\begin{table}[H]
\centering
\captionsetup{justification=centering}
\caption{Teste de diagnóstico do resíduo para o total de desocupados - modelos univariado e multivariado}
\label{tab:diagdesoc}
\scalebox{1}{
\renewcommand{\arraystretch}{0.7}
\begin{tabular}{lccc}
\toprule
& \multicolumn{3}{c}{\textbf{Modelo Univariado}} \\
\cmidrule(lr){2-4}
\textbf{Estrato geográfico} & Shapiro-Wilk & Ljung-Box & H \\
\midrule
01 - Belo Horizonte & 0,9769 & 0,6518 & 0,6333 \\
02 - Entorno e Colar Metropolitano de BH & 0,4622 & 0,0530* & 0,7115 \\
03 - Sul de Minas & 0,0917* & 0,0607* & 0,8793 \\
04 - Triângulo Mineiro & 0,0539* & 0,4577 & 0,9610 \\
05 - Zona da Mata & 0,6087 & 0,2385 & 0,6497 \\
06 - Norte de Minas & 0,7141 & 0,0848* & 0,7214 \\
07 - Vale do Rio Doce & 0,9441 & 0,8272 & 0,0997* \\
08 - Central & 0,7268 & 0,0928* & 0,1300 \\
\midrule
& \multicolumn{3}{c}{\textbf{Modelo Multivariado}} \\
\cmidrule(lr){2-4}
\textbf{Estrato geográfico} & Shapiro-Wilk & Ljung-Box & H \\
\midrule
01 - Belo Horizonte & 0,7846 & 0,1710 & 0,3265 \\
02 - Entorno e Colar Metropolitano de BH & 0,6429 & 0,0684* & 0,7184 \\
03 - Sul de Minas & 0,1143 & 0,0743* & 0,7924 \\
04 - Triângulo Mineiro & 0,7553 & 0,2309 & 0,3836 \\
05 - Zona da Mata & 0,1534 & 0,2716 & 0,8799 \\
06 - Norte de Minas & 0,3554 & 0,0010*** & 0,2680 \\
07 - Vale do Rio Doce & 0,0404** & 0,3737 & 0,1733 \\
08 - Central & 0,7102 & 0,7748 & 0,6216 \\
\bottomrule
\end{tabular}}
\fonte{Elaboração própria, com base nos dados da PNAD Contínua. Nota: * representa rejeição de \(H_0\) a 10\%, ** a 5\% e *** a 1\%. As oito primeiras observações foram descartadas, em razão do período de estabilização do filtro de Kalman.}
\end{table}
